\section{Introduction and the geometric determination theorem}
\label{sec:introduction}
A probability of a short collision record is an integrated geometric
quantity. It is not a marked length spectrum, and it need not retain the
orientation information present in a collision position. This distinction
is decisive for an inverse problem at a pair of reflecting contacts.
Reflection of both transverse coordinates preserves every unmarked
near-onset count. A signed first moment changes sign. We show that this
small, explicitly specified enrichment removes the evenness assumption
from two-contact recovery.

The central result is an all-degree inverse for two flights. The probability
coefficients recover successive even graph derivatives once the lower
ones are known. Signed endpoint first moments recover the intervening odd
derivatives. The two orientations give a nonsingular two-by-two block at
every degree. In particular, there is no exceptional locus at equal
curvatures and no generic nonvanishing assumption on a higher contact jet.
The observation used here is only a mean of the sum of two transverse
coordinates; it is not a measurement of the longitudinal boundary graph.

There is a complementary long-record invariant. In smooth dispersing
billiards the normalized mixed endpoint derivative has a relative product
limit on a nonshrinking collar. Residual-time integration produces two
symmetrized energy profiles and their mixed Volterra law. The complete
smooth construction, function-valued inverse, physical even-contact
specialization, and charged acquisition arguments are included in
\cite{Companion}. Section~\ref{sec:relative} states how they fit with the
present inverse, without identifying a count-only experiment with a
signed experiment.

The closest analytic forward comparison comes first: local hyperbolic
normal form, followed by projection to physical endpoints, already
explains an analytic relative product. The construction recalled in
\cite[Section 2]{DSKL} supplies that comparator. The directly geometric
smooth proof supplies common physical collars and differentiated
residual-time integration without assuming analytic normalizing charts.
Neither forward explanation alone recovers an asymmetric contact from
the scalar moments studied below. The global marked-length determination
in \cite{DSKL} uses a different information set and symmetry and genericity
hypotheses. Our data do not supply that spectrum; our conclusion concerns
the two participating boundaries in supplied contact frames.

\subsection{Setting and information}
Consider finitely many positively curved smooth obstacles modulo a fixed
full-rank lattice in $\R^2$, with disjoint closures and positive separation.
Let $A>0$ be the free cell area. Initial conditions have normalized phase
volume $\dd q\dd\vartheta/(2\pi A)$ and speed one. Fix an unobstructed
normal channel of length $g>0$ and label its contacts $0,1$.
Choose a single signed transverse axis at both contacts. In these frames
the two facing graphs have inward height functions
\begin{equation}\label{eq:graphs}
 \psi_b(y)=\frac{\kappa_b y^2}{2}
       +\sum_{r\ge3}\frac{q_{b,r}}{r!}y^r,\qquad \kappa_b>0,
       \quad b=0,1.
\end{equation}
For smooth graphs this formula denotes Taylor jets, not a convergent
expansion. The gap, labels, axis orientation, and separate curvatures are
supplied. Area may be unknown. No congruence, reflection symmetry, or
finite-horizon assumption is made.

Fix small facing patches before observation. Let $E_{2,b}(d)$ be the event
that the time window $2g+d$ contains exactly the three successive impacts
$b,1-b,b$ in these patches. Let $u,v$ be the first and last transverse
coordinates. Put
\begin{equation}\label{eq:rawdata}
 P_b(d)=\Prob(E_{2,b}(d)),\qquad
 R_b(d)=\E\bigl[(u+v)\one_{E_{2,b}(d)}\bigr].
\end{equation}
The random variable in the second formula is set to zero on failure.
In particular this is an unconditional moment, not a conditional mean
with an uncharged selection step. Only the two scalar functions per
orientation are the exact data for the inverse theorem.

Write
\begin{equation}\label{eq:leading}
 c_b=1+g\kappa_b,\quad z=c_0c_1-1>0,\quad
 \gamma=\arcosh\sqrt{c_0c_1},\quad
 \alpha=\frac1{2A\sinh(2\gamma)}.
\end{equation}
The normalized germs and their coefficients are
\begin{equation}\label{eq:coefficients}
 G_b(d)=\frac{P_b(d)}{\alpha d^2}=1+\sum_{k\ge1}\xi_{b,k}d^k,
 \qquad
 J_b(d)=\frac{R_b(d)}{\alpha d^2}=\sum_{k\ge1}\eta_{b,k}d^k.
\end{equation}
Their smooth extension to $d=0$ is part of the proof.

\begin{theorem}[Asymmetric two-contact determination]\label{thm:main}
For the setting above the following assertions hold.

\smallskip\noindent
\textup{(i)} The common leading coefficient of the raw probability germs
is $\alpha=\lim_{d\downarrow0}P_b(d)/d^2$. It determines $A$ from the
supplied leading geometry.

\smallskip\noindent
\textup{(ii)} For every fixed $K\ge3$, set
$n_e=\lfloor K/2\rfloor-1$ and $n_o=\lfloor(K-1)/2\rfloor$.
The coefficient vector
\[
 (\xi_{b,k})_{b=0,1;\,1\le k\le n_e},\qquad
 (\eta_{b,k})_{b=0,1;\,1\le k\le n_o}
\]
is an analytic, block-triangular coordinate system for
$(q_{b,r})_{b=0,1;\,3\le r\le K}$.
The inverse is explicit by successive two-by-two solves, with blocks
\eqref{eq:evenblock} and \eqref{eq:oddblock}. At each fixed order it is
locally bi-Lipschitz on compact finite-jet neighborhoods and compact
positive leading-geometry boxes, including $\kappa_0=\kappa_1$.

\smallskip\noindent
\textup{(iii)} If both contacts are real analytic, equality of the four
raw germs in \eqref{eq:rawdata}, at the same supplied leading geometry,
determines $A$ and both contact graph germs. For connected analytic
closed obstacle boundaries it determines the two participating boundary
images in their supplied frames.

\smallskip\noindent
\textup{(iv)} At every fixed $K$ there are analytic periodic physical
families with independent graph jets of every degree from $3$ to $K$,
either with area fixed or with area as an additional coordinate. On
these supplied families, respectively $2K-4$ or $2K-3$ positive-window
scalar means are local bi-Lipschitz coordinates. In the fixed-area
case the selected leading count and endpoint-metric hierarchy is fixed
at every flight number.
\end{theorem}

Parts (i)--(iii) are proved in Sections~\ref{sec:local}--\ref{sec:inverse};
part (iv) is proved in Section~\ref{sec:physical}. All orders are fixed
before neighborhoods and constants are chosen. Equality of all smooth
jets is not asserted to identify a smooth germ. Recovery from exact
germs is not a finite noisy measurement; the positive-window assertion
is a separate theorem on a specified physical family.

The scalar count in part (iv) is the dimension of a differentiable
locally bi-Lipschitz mean coordinate map. It is not a lower bound over
all experiments. The statistically compressed experiment is stated
explicitly in Theorem~\ref{thm:experiment}; its lower risk bound is not
claimed for the uncompressed endpoint record.
